\documentclass[10pt,letterpaper]{article}
\usepackage[utf8]{inputenc}
\usepackage[intlimits]{amsmath}
\usepackage{amsfonts}
\usepackage{amssymb}
\usepackage{ragged2e}
\usepackage[letterpaper, margin=0.5in]{geometry}
\usepackage{graphicx}
\usepackage{cancel}
\usepackage{mathtools}
\usepackage{tabularx}
\usepackage{arydshln}
\usepackage{tensor}
\usepackage{array}
\usepackage{xcolor}
\usepackage[boxed]{algorithm}
\usepackage[noend]{algpseudocode}
\usepackage{listings}
\usepackage{textcomp}
% \usepackage[pdf,tmpdir,singlefile]{graphviz}
\usepackage{mathrsfs}
\usepackage{bbm}
\usepackage{tikz}
\usepackage{tikz-cd}
\usepackage{enumitem}
\usepackage{arydshln}
\usepackage{relsize}
\usepackage{multicol}
\usepackage{scalerel}
\usepackage{upgreek}
\usepackage{wasysym}

\usetikzlibrary{bayesnet}
\setlist{noitemsep}

%%%%%%%%%%%%%%%%%%%%%%%%%%%%%
% Formatting commands
%%%%%%%%%%%%%%%%%%%%%%%%%%%%%
\newcommand{\n}{\hfill\break}
\newcommand{\nn}{\vspace{0.5\baselineskip}\n}
\newcommand{\up}{\vspace{-\baselineskip}}
\newcommand{\hangblock}[2]{\par\noindent\settowidth{\hangindent}{\textbf{#1: }}\textbf{#1: }\nolinebreak#2}
\newcommand{\lemma}[1]{\hangblock{Lemma}{#1}}
\newcommand{\defn}[1]{\hangblock{Defn}{#1}}
\newcommand{\thm}[1]{\hangblock{Thm}{#1}}
\newcommand{\cor}[1]{\hangblock{Cor}{#1}}
\newcommand{\prop}[1]{\hangblock{Prop}{#1}}
\newcommand{\ex}[1]{\hangblock{Ex}{#1}}
\newcommand{\exer}[1]{\hangblock{Exer}{#1}}
\newcommand{\fact}[1]{\hangblock{Fact}{#1}}
\newcommand{\remark}[1]{\hangblock{Remark}{#1}}
\newcommand{\proven}{\;$\square$\n}
\newcommand{\problem}[1]{\par\noindent{#1}\n}
\newcommand{\problempart}[2]{\par\noindent\indent{}\settowidth{\hangindent}{\textbf{(#1)} \indent{}}\textbf{(#1) }\nolinebreak#2\n}
\newcommand{\ptxt}[1]{\textrm{\textnormal{#1}}}
\newcommand{\inlineeq}[1]{\centerline{$\displaystyle #1$}}
\newcommand{\pageline}{\up\par\noindent\rule{\textwidth}{0.1pt}}

%%%%%%%%%%%%%%%%%%%%%%%%%%%%%
% Math commands
%%%%%%%%%%%%%%%%%%%%%%%%%%%%%
% Set Theory
\newcommand{\card}[1]{\left|#1\right|}
\newcommand{\set}[1]{\left\{#1\right\}}
\newcommand{\setmid}{\;\middle|\;}
\newcommand{\ps}[1]{\mathcal{P}\left(#1\right)}
\newcommand{\pfinite}[1]{\mathcal{P}^{\ptxt{finite}}\left(#1\right)}
\newcommand{\naturals}{\mathbb{N}}
\newcommand{\N}{\naturals}
\newcommand{\integers}{\mathbb{Z}}
\newcommand{\Z}{\integers}
\newcommand{\rationals}{\mathbb{Q}}
\newcommand{\Q}{\rationals}
\newcommand{\reals}{\mathbb{R}}
\newcommand{\R}{\reals}
\newcommand{\complex}{\mathbb{C}}
\newcommand{\C}{\complex}
\newcommand{\halfPlane}{\mathbb{H}}
\let\HSym\H
\let\H\relax
\newcommand{\H}{\halfPlane}
\newcommand{\comp}{^{\complement}}
\DeclareMathOperator{\Hom}{Hom}
\newcommand{\Ind}{\mathbbm{1}}
\newcommand{\cut}{\setminus}
\DeclareMathOperator{\elem}{elem}

% Differentiable Manifolds
\newcommand{\RP}{\mathbb{RP}}
\newcommand{\CP}{\mathbb{RP}}
\newcommand{\osubset}{\overset{\mathclap{\scalebox{0.5}{\ptxt{open}}}}{\subset}}
\newcommand{\osubseteq}{\overset{\mathclap{\scalebox{0.5}{\ptxt{open}}}}{\subseteq}}
\newcommand{\osupset}{\overset{\mathclap{\scalebox{0.5}{\ptxt{open}}}}{\supset}}
\newcommand{\osupseteq}{\overset{\mathclap{\scalebox{0.5}{\ptxt{open}}}}{\supseteq}}
\newcommand{\pdat}[3]{\left.\pd{#1}{#2}\right|_{#3}}
\DeclareMathOperator{\so}{so}
\DeclareMathOperator{\codim}{codim}
\DeclareMathOperator{\Diff}{Diff}
\let\dSym\d
\let\d\relax
\newcommand{\d}{\partial}
\DeclareMathOperator{\gl}{gl}

% Graph Theory
\let\deg\relax
\DeclareMathOperator{\deg}{deg}
\newcommand{\degp}{\ptxt{deg}^{+}}
\newcommand{\degn}{\ptxt{deg}^{-}}
\newcommand{\precdot}{\mathrel{\ooalign{$\prec$\cr\hidewidth\hbox{$\cdot\mkern0.5mu$}\cr}}}
\newcommand{\succdot}{\mathrel{\ooalign{$\cdot\mkern0.5mu$\cr\hidewidth\hbox{$\succ$}\cr\phantom{$\succ$}}}}
\DeclareMathOperator{\cl}{cl}
\DeclareMathOperator{\affdim}{affdim}

% Probability
\newcommand{\parSymbol}{\P}
\newcommand{\Prob}{\mathbb{P}}
\renewcommand{\P}{\Prob}
\newcommand{\Avg}{\mathbb{E}}
\newcommand{\E}{\Avg}
\DeclareMathOperator{\Var}{Var}
\DeclareMathOperator{\cov}{cov}
\DeclareMathOperator{\Unif}{Unif}
\DeclareMathOperator{\Binom}{Binom}
\newcommand{\CI}{\mathrel{\text{\scalebox{1.07}{$\perp\mkern-10mu\perp$}}}}

% Standard Math
\newcommand{\inv}{^{-1}}
\newcommand{\abs}[1]{\left|#1\right|}
\newcommand{\ceil}[1]{\left\lceil{}#1\right\rceil{}}
\newcommand{\floor}[1]{\left\lfloor{}#1\right\rfloor{}}
\newcommand{\conj}[1]{\overline{#1}}
\newcommand{\of}{\circ}
\newcommand{\tri}{\triangle}
\newcommand{\inj}{\hookrightarrow}
\newcommand{\surj}{\twoheadrightarrow}
\newcommand{\ndiv}{\nmid}
\renewcommand{\epsilon}{\varepsilon}
\newcommand{\divides}{\mid}
\newcommand{\ndivides}{\nmid}
\DeclareMathOperator{\lcm}{lcm}
\DeclareMathOperator{\sgn}{sgn}
\newcommand{\map}[4]{\!\!\!\begin{array}[t]{rcl}#1 & \!\!\!\!\to & \!\!\!\!#2\\ {}#3 & \!\!\!\!\mapsto & \!\!\!\!#4\end{array}}
\newcommand{\bigsum}[2]{\smashoperator[lr]{\sum_{\scalebox{#1}{$#2$}}}}
\DeclareMathOperator{\gcf}{gcf}
\newcommand{\restr}[2]{\left.#1\right|_{#2}}

% Linear Algebra
\newcommand{\Id}{\textrm{\textnormal{Id}}}
\newcommand{\im}{\textrm{\textnormal{im}}}
\newcommand{\norm}[1]{\abs{\abs{#1}}}
\newcommand{\tpose}{^{T}\!}
\newcommand{\iprod}[1]{\left<#1\right>}
\newcommand{\giprod}{\iprod{\;\,,\;}}
\DeclareMathOperator{\tr}{tr}
\DeclareMathOperator{\trace}{tr}
\newcommand{\chgBasMat}[3]{\!\!\tensor*[_{#1}]{\left[#2\right]}{_{#3}}}
\newcommand{\vecBas}[2]{\tensor*[]{\left[#1\right]}{_{#2}}}
\DeclareMathOperator{\GL}{GL}
\DeclareMathOperator{\Mat}{Mat}
\DeclareMathOperator{\vspan}{span}
\DeclareMathOperator{\rank}{rank}
\newcommand{\V}[1]{\vec{#1}}
\DeclareMathOperator{\proj}{proj}
\DeclareMathOperator{\compProj}{comp}
\DeclareMathOperator{\row}{row}
\newcommand{\smallPMatrix}[1]{\paren{\begin{smallmatrix}#1\end{smallmatrix}}}
\newcommand{\smallBMatrix}[1]{\brack{\begin{smallmatrix}#1\end{smallmatrix}}}
\newcommand{\pmat}[1]{\begin{pmatrix}#1\end{pmatrix}}
\newcommand{\bmat}[1]{\begin{bmatrix}#1\end{bmatrix}}
\newcommand{\dual}{^{*}}
\newcommand{\pinv}{^{\dagger}}
\newcommand{\horizontalMatrixLine}{\ptxt{\rotatebox[origin=c]{-90}{$|$}}}
\DeclareMathOperator{\range}{range}
\DeclareMathOperator{\Symm}{Symm}
\DeclareMathOperator{\SU}{SU}

% Multilinear Algebra
\let\Lsym\L
\let\L\relax
\DeclareMathOperator{\L}{\mathscr{L}}
\DeclareMathOperator{\A}{\mathcal{A}}
\DeclareMathOperator{\Alt}{Alt}
\DeclareMathOperator{\Sym}{Sym}
\newcommand{\ot}{\otimes}
\newcommand{\ox}{\otimes}
\DeclareMathOperator{\asc}{asc}
\DeclareMathOperator{\asSet}{set}
\DeclareMathOperator{\sort}{sort}
\DeclareMathOperator{\ringA}{\mathring{A}}
\DeclareMathOperator{\Sh}{Sh}

% Topology
\newcommand{\closure}[1]{\overline{#1}}
\newcommand{\uball}{\mathcal{U}}
\DeclareMathOperator{\Int}{Int}
\DeclareMathOperator{\Ext}{Ext}
\DeclareMathOperator{\Bd}{Bd}
\DeclareMathOperator{\rInt}{rInt}
\DeclareMathOperator{\ch}{ch}
\DeclareMathOperator{\ah}{ah}
\newcommand{\LargerTau}{\mathlarger{\mathlarger{\mathlarger{\mathlarger{\tau}}}}}
\newcommand{\Tau}{\mathcal{T}}

% Analysis
\DeclareMathOperator{\Graph}{Graph}
\DeclareMathOperator{\epi}{epi}
\DeclareMathOperator{\hypo}{hypo}
\DeclareMathOperator{\supp}{supp}
\newcommand{\lint}[2]{\underset{#1}{\overset{#2}{{\color{black}\underline{{\color{white}\overline{{\color{black}\int}}\color{black}}}}}}}
\newcommand{\uint}[2]{\underset{#1}{\overset{#2}{{\color{white}\underline{{\color{black}\overline{{\color{black}\int}}\color{black}}}}}}}
\newcommand{\alignint}[2]{\underset{#1}{\overset{#2}{{\color{white}\underline{{\color{white}\overline{{\color{black}\int}}\color{black}}}}}}}
\newcommand{\extint}{\ptxt{ext}\int}
\newcommand{\extalignint}[2]{\ptxt{ext}\alignint{#1}{#2}}
\newcommand{\conv}{\ast}
\newcommand{\pd}[2]{\frac{\partial{}#1}{\partial{}#2}}
\newcommand{\del}{\nabla}
\DeclareMathOperator{\grad}{grad}
\DeclareMathOperator{\curl}{curl}
\let\div\relax
\DeclareMathOperator{\div}{div}
\DeclareMathOperator{\vol}{vol}

% Complex Analysis
\let\Re\relax
\DeclareMathOperator{\Re}{Re}
\let\Im\relax
\DeclareMathOperator{\Im}{Im}
\DeclareMathOperator{\Res}{Res}

% Abstract Algebra
\DeclareMathOperator{\ord}{ord}
\newcommand{\generated}[1]{\left<#1\right>}
\newcommand{\cycle}[1]{\smallPMatrix{#1}}
\newcommand{\id}{\ptxt{id}}
\newcommand{\iso}{\cong}
\DeclareMathOperator{\Aut}{Aut}
\DeclareMathOperator{\SL}{SL}
\DeclareMathOperator{\op}{op}
\newcommand{\isom}[4]{\!\!\!\begin{array}[t]{rcl}#1 & \!\!\!\!\overset{\sim}{\to} & \!\!\!\!#2\\ #3 & \!\!\!\!\mapsto & \!\!\!\!#4\end{array}}
\newcommand{\F}{\mathbb{F}}
\newcommand{\acts}{\;\reflectbox{\rotatebox[origin=c]{-90}{$\circlearrowright$}}\;}
\newcommand{\disjunion}{\mathrel{\text{\scalebox{1.07}{$\perp\mkern-10mu\perp$}}}}
\DeclareMathOperator{\SO}{SO}
\DeclareMathOperator{\stab}{stab}
\DeclareMathOperator{\nullity}{nullity}
\DeclareMathOperator{\Perm}{Perm}
\DeclareMathOperator{\nsubgp}{\vartriangleleft}
\DeclareMathOperator{\notnsubgp}{\ntriangleleft}
\newcommand{\presentation}[2]{\left<#1\;\middle|\;#2\right>}
\DeclareMathOperator{\Char}{char}
\DeclareMathOperator{\fchar}{char}
\DeclareMathOperator{\triv}{triv}
\DeclareMathOperator{\reg}{reg}
\DeclareMathOperator{\std}{std}
\DeclareMathOperator{\Func}{Func}
\DeclareMathOperator{\End}{End}

% Convex Optimization
\let\sectionSymbol\S
\let\S\relax
\newcommand{\S}{\mathbb{S}}
\DeclareMathOperator{\dist}{dist}
\DeclareMathOperator{\dom}{dom}
\DeclareMathOperator{\diag}{diag}
\DeclareMathOperator{\ones}{\mathbbm{1}}
\newcommand{\minimizeOver}[3]{\begin{array}{rl}\underset{#1}{\ptxt{minimize}} & #2\\ \ptxt{subject to} & #3\end{array}}
\newcommand{\maximizeOver}[3]{\begin{array}{rl}\underset{#1}{\ptxt{maximize}} & #2\\ \ptxt{subject to} & #3\end{array}}
\newcommand{\minimizationProblem}[2]{\minimizeOver{}{#1}{#2}}
\newcommand{\maximizationProblem}[2]{\maximizeOver{}{#1}{#2}}
\newcommand{\minimizeOverUnconstrained}[2]{\begin{array}{rl}\underset{#1}{\ptxt{minimize}} & #2\end{array}}
\newcommand{\maximizeOverUnconstrained}[2]{\begin{array}{rl}\underset{#1}{\ptxt{maximize}} & #2\end{array}}
\newcommand{\minimizationUnconstrained}[1]{\minimizeOverUnconstrained{}{#1}}
\newcommand{\maximizationUnconstrained}[1]{\maximizeOverUnconstrained{}{#1}}
\DeclareMathOperator{\argmin}{argmin}
\DeclareMathOperator{\argmax}{argmax}

% Proofs
\newcommand{\st}{s.t.}
\newcommand{\unique}{!}
\newcommand{\iffdef}{\overset{\ptxt{def}}{\Leftrightarrow}}
\newcommand{\eqVertical}{\rotatebox[origin=c]{90}{=}}
\newcommand{\mapsfrom}{\mathrel{\reflectbox{\ensuremath{\mapsto}}}}
\newcommand{\mapsdown}{\rotatebox[origin=c]{-90}{$\mapsto$}\mkern2mu}
\newcommand{\mapsup}{\rotatebox[origin=c]{90}{$\mapsto$}\mkern2mu}
\newcommand{\from}{\!\mathrel{\reflectbox{\ensuremath{\to}}}}
\newcommand{\labeledeq}[1]{\overset{\mathclap{\ptxt{#1}}}{=}}
\newcommand{\eqdef}{\labeledeq{def}}

% Brackets
\newcommand{\paren}[1]{\left(#1\right)}
\renewcommand{\brack}[1]{\left[#1\right]}
\renewcommand{\brace}[1]{\left\{#1\right\}}
\newcommand{\ang}[1]{\left<#1\right>}

% Algorithms
\algrenewcommand{\algorithmiccomment}[1]{\hskip 1em \texttt{// #1}}
\algrenewcommand\algorithmicrequire{\textbf{Input:}}
\algrenewcommand\algorithmicensure{\textbf{Output:}}
\newcommand{\algP}{\ptxt{\textbf{P}}}
\newcommand{\algNP}{\ptxt{\textbf{NP}}}
\newcommand{\algNPC}{\ptxt{\textbf{NP-Complete}}}
\newcommand{\algNPH}{\ptxt{\textbf{NP-Hard}}}
\newcommand{\algEXP}{\ptxt{\textbf{EXP}}}
\DeclareMathOperator{\fl}{fl}

%%%%%%%%%%%%%%%%%%%%%%%%%%%%%
% Other commands
%%%%%%%%%%%%%%%%%%%%%%%%%%%%%
\newcommand{\flag}[1]{\textbf{\textcolor{red}{#1}}}
\let\uSym\u
\let\u\relax
\newcommand{\u}[1]{\underline{#1}}
\let\bSym\b
\let\b\relax
\newcommand{\b}[1]{\textbf{#1}}
\let\iSym\i
\let\i\relax
\newcommand{\i}[1]{\textit{#1}}
\let\scSym\sc
\let\sc\relax
\newcommand{\sc}[1]{\textsc{#1}}

%%%%%%%%%%%%%%%%%%%%%%%%%%%%%%%%%%%%%%%
% Make l's curvy in math environments %
%%%%%%%%%%%%%%%%%%%%%%%%%%%%%%%%%%%%%%%
\mathcode`l="8000
\begingroup
\makeatletter
\lccode`\~=`\l
\DeclareMathSymbol{\lsb@l}{\mathalpha}{letters}{`l}
\lowercase{\gdef~{\ifnum\the\mathgroup=\m@ne \ell \else \lsb@l \fi}}%
\endgroup

%%%%%%%%%%%%%%%%%%%%%%%%%
% Fix \vdots and \ddots %
%%%%%%%%%%%%%%%%%%%%%%%%%
\usepackage{letltxmacro}
\LetLtxMacro\orgvdots\vdots
\LetLtxMacro\orgddots\ddots

\makeatletter
\DeclareRobustCommand\vdots{%
	\mathpalette\@vdots{}%
}
\newcommand*{\@vdots}[2]{%
	% #1: math style
	% #2: unused
	\sbox0{$#1\cdotp\cdotp\cdotp\m@th$}%
	\sbox2{$#1.\m@th$}%
	\vbox{%
		\dimen@=\wd0 %
		\advance\dimen@ -3\ht2 %
		\kern.5\dimen@
		% remove side bearings
		\dimen@=\wd2 %
		\advance\dimen@ -\ht2 %
		\dimen2=\wd0 %
		\advance\dimen2 -\dimen@
		\vbox to \dimen2{%
			\offinterlineskip
			\copy2 \vfill\copy2 \vfill\copy2 %
		}%
	}%
}
\DeclareRobustCommand\ddots{%
	\mathinner{%
		\mathpalette\@ddots{}%
		\mkern\thinmuskip
	}%
}
\newcommand*{\@ddots}[2]{%
	% #1: math style
	% #2: unused
	\sbox0{$#1\cdotp\cdotp\cdotp\m@th$}%
	\sbox2{$#1.\m@th$}%
	\vbox{%
		\dimen@=\wd0 %
		\advance\dimen@ -3\ht2 %
		\kern.5\dimen@
		% remove side bearings
		\dimen@=\wd2 %
		\advance\dimen@ -\ht2 %
		\dimen2=\wd0 %
		\advance\dimen2 -\dimen@
		\vbox to \dimen2{%
			\offinterlineskip
			\hbox{$#1\mathpunct{.}\m@th$}%
			\vfill
			\hbox{$#1\mathpunct{\kern\wd2}\mathpunct{.}\m@th$}%
			\vfill
			\hbox{$#1\mathpunct{\kern\wd2}\mathpunct{\kern\wd2}\mathpunct{.}\m@th$}%
		}%
	}%
}
\makeatother

\newcommand{\B}{
	\begin{tikzpicture}
	\filldraw [fill=red, draw=black] (0, 0) rectangle (0.37, 0.45);
	\draw [line width=0.5mm, white ] (0.1,0.08) -- (0.1,0.38);
	\draw[line width=0.5mm, white ] (0.1, 0.35) .. controls (0.2, 0.35) and (0.4, 0.2625) .. (0.1, 0.225);
	\draw[line width=0.5mm, white ] (0.1, 0.225) .. controls (0.2, 0.225) and (0.4, 0.1625) .. (0.1, 0.1);
	\end{tikzpicture}
}

% Allow custom symbols as arrows in tikz-cd
\tikzset{
  symbol/.style={
    draw=none,
    every to/.append style={
      edge node={node [sloped, allow upside down, auto=false]{$#1$}}}
  }
}
% Example: \arrow[r,symbol=\cong]
% https://tex.stackexchange.com/questions/394154/how-to-include-inclusion-subgroup-relationship-in-tikz-cd-diagram

\author{Thomas Cohn}
\title{Math 591 Lecture 32}
\date{11/16/20} % Can also use \today

\begin{document}
\maketitle
\setlength\RaggedRightParindent{\parindent}
\RaggedRight

\subsection*{The Exterior Differential}

\thm{
	Let $M$ be a manifold. Then there exists a unique operator $d$
	\[
		\Omega^{0}\overset{d}{\xrightarrow{\hspace{0.5cm}}}\Omega^{1}\overset{d}{\xrightarrow{\hspace{0.5cm}}}\Omega^{2}\overset{d}{\xrightarrow{\hspace{0.5cm}}}\cdots\overset{d}{\xrightarrow{\hspace{0.5cm}}}\Omega^{k-1}\overset{d}{\xrightarrow{\hspace{0.5cm}}}\Omega^{k}
	\]
	\st{}
	\begin{enumerate}[itemsep=0pt, topsep=0pt, leftmargin=4\parindent, label=(\arabic*)]
		\item $d$ is $\R$-linear.
		\item $d:\Omega^{0}\to\Omega^{1}$ is the usual differential of a function.
		\item $d^{2}=0$.
		\item The anti-derivation property: $d(\alpha\wedge\beta)=d\alpha\wedge\beta+(-1)^{\deg\alpha}\alpha\wedge{}d\beta$. (This is also known as the super-symmetric version of Leibniz' rule.)
		\item If $F:M\to{}N$ is smooth, then $\forall\alpha\in\Omega^{k}N$, $d(F^{*}\alpha)=F^{*}(d\alpha)$. (This property is called ``Naturality''.)
	\end{enumerate}\up\n
}

\par\noindent
Note: In theory, the $d$ is a different operator for each $\Omega^{i}$. But it would be inconvenient to call them different names.\n

\par\noindent
Proof: We'll begin by proving this for $U\osubseteq\R^{n}$, and then use coordinates. This works because the anti-derivation property implies locality.\n

\par\noindent
For $\R^{n}$: If $\displaystyle\alpha=\sum_{I}a_{I}dx^{I}\in\Omega^{k}(\R^{n})$, then the properties imply that we must have $\displaystyle{}d\alpha=\sum_{I}\sum_{j=1}^{n}\pd{a_{I}}{x^{j}}dx^{j}\wedge{}dx^{I}$.\n
Use this as our definition of $d$ on $\R^{n}$. Now, we must check the $5$ properties. $1$ and $4$ follow directly from its definition. By linearlity, we can check this solely on ``monomials''.\n

\par\noindent
\b{Anti-derivation:} Let $\alpha=a\,dx^{I}$, $\beta=b\,dx^{J}$, $a,b\in{}C^{\infty}$. Then
\begin{align*}
	d(\alpha\wedge\beta) & =d(a\,dx^{I}\wedge{}b\,dx^{J})\\
	& =d(ab\,dx^{I}\wedge{}dx^{J})\\
	& =d(ab\,dx^{IJ})\\
	& =\sum_{j=1}^{n}\pd{(ab)}{x^{j}}\,dx^{j}\wedge{}dx^{IJ}\\
	& =\sum_{j=1}^{n}\brack{b\pd{a}{x^{j}}+a\pd{b}{x^{j}}}\,dx^{j}\wedge{}dx^{IJ}\\
	& =\sum_{j=1}^{n}\pd{a}{x^{j}}\,dx^{j}\wedge{}dx^{I}\wedge(b\,dx^{J})+a\sum_{j=1}^{n}\paren{\pd{b}{x^{j}}\,dx^{j}\wedge{}dx^{I}}\wedge{}dx^{J}\\
	& =\sum_{j=1}^{n}\pd{a}{x^{j}}\,dx^{j}\wedge{}dx^{I}\wedge(b\,dx^{J})+(-1)^{k}a\,dx^{I}\sum_{j=1}^{n}\pd{b}{x^{j}}\,dx^{j}\wedge{}dx^{I}\\
	& =d\alpha\wedge\beta+(-1)^{k}\alpha\wedge{}d\beta
\end{align*}
\n

\newpage
\par\noindent
\b{$\mathbf{d^{2}=0}$:}
\begin{align*}
	d\brack{d(a\,dx^{I})} & =d\brack{\sum_{j=1}^{n}\pd{a}{x^{j}}\,dx^{j}\wedge{}dx^{I}}\\
	& \labeledeq{(1)}\sum_{j=1}^{n}d\paren{\pd{a}{x^{j}}}\wedge{}dx^{j}\wedge{}dx^{I}\\
	& =\underbrace{\paren{\sum_{j=1}^{n}\sum_{l=1}^{n}\frac{\d^{2}a}{\d{}x^{l}\d{}x^{j}}\,dx^{l}\wedge{}dx^{j}}}_{\mathclap{\ptxt{Observe: Terms with $j=l$ vanish.}}}\wedge{}dx^{I}\\
	& =\paren{\sum_{1\le{}j<l\le{}n}\frac{\d^{2}a}{\d{}x^{l}\d{}x^{j}}\,dx^{l}\wedge{}dx^{j}+\frac{\d^{2}a}{\d{}x^{j}\d{}x^{l}}\,dx^{j}\wedge{}dx^{l}}\wedge{}dx^{I}\\
	& =\paren{\sum_{1\le{}j<l\le{}n}\paren{\frac{\d^{2}a}{\d{}x^{l}\d{}x^{j}}-\frac{\d^{2}}{\d{}x^{j}\d{}x^{l}}}\,dx^{l}\wedge{}dx^{j}}\wedge{}dx^{I}\\
	& \labeledeq{(2)}0
\end{align*}
with (1) because $dx^{j}\wedge{}dx^{I}$ is constant, so $d(dx^{j}\wedge{}dx^{I})=0$, and (2) by Clairaut's Theorem, which tells us that\break $\frac{\d^{2}a}{\d{}x^{l}\d{}x^{j}}=\frac{\d^{2}a}{\d{}x^{j}\d{}x^{l}}$, so their difference is $0$.\n

\par\noindent
\b{Naturality:} Say $F:V\to{}U$ is $C^{\infty}$, $V\subseteq\R^{m}$, $U\subseteq\R^{n}$. Let $\alpha\in\Omega^{k}(U)$ of the form $\alpha=a\,dx^{I}$. We need to prove that $d(F^{*}\alpha)=F^{*}(d\alpha)$. Say $F=(F^{1},\ldots,F^{n})$ with $F^{j}:V\to\R$. We proved that
\[
	F^{*}(dx^{i_{1}})\wedge\cdots\wedge{}F^{*}(dx^{i_{k}})=F^{*}(dx^{I})=dF^{i_{1}}\wedge\cdots\wedge{}dF^{i_{k}}\eqdef{}dF^{I}
\]
So $F^{*}\alpha=(\alpha\of{}F)F^{*}(dx^{I})=(\alpha\of{}F)dF^{I}$, and thus,
\[
	d(F^{*}\alpha)=d(\alpha\of{}F)\wedge{}dF^{I}+(\alpha\of{}F)\underbrace{d(dF^{I})}_{=0}=d(\alpha\of{}F)\wedge{}dF^{I}
\]
\[
	F^{*}(d\alpha)=F^{*}\brack{\sum_{j=1}^{n}\pd{a}{x^{j}}\,dx^{j}\wedge{}dx^{I}}=\sum_{j=1}^{n}F^{*}\paren{\pd{a}{x^{j}}}\,dx^{j}\wedge{}dx^{I}=\sum_{j=1}^{n}\paren{\pd{a}{x^{j}}\of{}F}\underbrace{\paren{F^{*}(dx^{j})}}_{dF^{j}}\wedge{}dF^{I}=d(\alpha\of{}F)\wedge{}dF^{I}
\]
So we conclude that $d(F^{*}\alpha)=F^{*}(d\alpha)$.\n

\par\noindent
Now, we use coordinates to do it on a manifold.\n

\defn{
	Given $\alpha\in\Omega^{k}(M)$, $\phi:U\to\R^{n}$ a coordinate system, define $d\alpha\in\Omega^{k+1}(U)$ by $\phi^{*}\brack{d(\phi\inv)^{*}\alpha}$, where $d$ is the Euclidean version, and $(\phi\inv)^{*}\alpha\in\Omega^{k}(\phi(U))$.\n
}

\par\noindent
Now, we need to prove that $d\alpha$ on $U$ is independent of choice of coordinates. Suppose $\psi:V\to\R^{n}$ is another coordinate chart (WOLOG $U\cap{}V\ne\emptyset$). Call $U$ the intersection of their domains. Write $\psi=F\of\phi$
\[
	\begin{tikzcd}
		&[-15pt] U \arrow[dl, "\phi", swap] \arrow[dr, "\psi"] &[-15pt]\\
		\phi(U) \arrow[rr, "F", swap] & & \psi(U)
	\end{tikzcd}
\]
Then
\begin{align*}
	\psi^{*}\brack{d(\psi\inv)^{*}\alpha} & =\phi^{*}\of{}F^{*}\brack{d(\phi\inv\of{}F\inv)^{*}\alpha}\\
	& =\phi^{*}\of{}F^{*}\brack{d((F\inv)^{*}\of(\phi\inv)^{*})\alpha}\\
	& =\phi^{*}\of{}F^{*}\of(F\inv)^{*}\brack{d(\phi\inv)^{*}\alpha}\\
	& =\phi^{*}\brack{d(\phi\inv)^{*}\alpha}
\end{align*}
\n

\par\noindent
As for showing uniqueness on manifolds, use that anti-derivation implies locality, as with vector fields.\proven

\ex{
	$\R^{3}$
	\begin{align*}
		d(f\,dx+g\,dy+h\,dz) & =df\wedge{}dx+dg\wedge{}dy+dh\wedge{}dz\\
		& =(\d_{x}f\,dx+\d_{y}f\,dy+\d_{z}f\,dz)+\cdots\\
		& =\d_{y}f\,dy\wedge{}dx+\d_{z}f\,dz\wedge{}dx+\cdots\\
		& =(\d_{x}g-\d_{y}f)\,dx\wedge{}dy+(\d_{x}h-\d_{z}f)\,dx\wedge{}dz+\cdots
	\end{align*}
	We get the components of the curl of the vector field $<f,g,h>$.\n
}

\defn{
	$\alpha\in\Omega^{k}(M)$ is \u{closed} if $d\alpha=0$, and \u{exact} if $\exists\beta\in\Omega^{k-1}(M)$ such that $\alpha=d\beta$.\n
}

\par\noindent
Note: $d^{2}=0$ means exact implies closed.\n

\defn{
	The \u{$k$th DeRahm Cohomology} $H^{k}(M)=\ker(d:\Omega^{k}\to\Omega^{k+1})/\im(d:\Omega^{k-1}\to\Omega^{k})$ (with the vector space quotient).\n
}

\ex{
	If $M=S^{1}$, $d\theta$ dual to $\pd{}{\theta}$ is closed. (Any top-degree form is closed.)\n
	Claim: $d\theta$ is not exact, and in fact, $H^{1}(S^{1})=\R$.\n
}

\end{document}